\documentclass[11pt]{article}

\usepackage[margin=1in]{geometry}
\usepackage{amsmath}
\usepackage{booktabs}
\usepackage{longtable}
\usepackage{array}
\usepackage{hyperref}

\hypersetup{colorlinks=true, linkcolor=black, urlcolor=blue, citecolor=black}

\title{Predicting Tommy Award Winners with Machine Learning:\\
From Boston Celtics Game Logs to League-Wide Player Profiles}
\author{Teddy Taussig}
\date{\today}

\begin{document}
\maketitle

\begin{abstract}
This project develops a machine learning model to predict Tommy Award winners among Boston Celtics players and then extend that model to the entire NBA. The Tommy Award is given to one Celtics player after every game for hustle and energy. The analysis uses ten seasons of award data (2016-17 to 2025-26) provided by Chris Forsberg of NBC Sports Boston, plus traditional statistics (points, rebounds, assists) from the NBA API. Models tested include Ridge, Lasso, and Elastic-net Logistic Regression, Decision Trees, Random Forest, and XGBoost Classifier. Hyperparameters are tuned with Optuna. The best-performing model was Random Forest, achieving a 78.4\% top-3 accuracy. This means the actual winner was among the model's top 3 predicted players 78\% of the time.
\end{abstract}

\section{Introduction}

\subsection{Motivation}
The Motivation for our research subject is the fact that we love the Boston Celtics. For the past three seasons, we have watched every game. Tommy Heinsohn passed away in 2020, yet the current broadcasters, Brian Scalabrine and Drew Carter, carry on his legacy. During the game, when a Celtics player makes an impressive hustle play, they’ll say, “That's a Tommy Point”. Or, when many players on the Celtics are playing well, Scalebrine will say something like, “I don’t know Drew, the Tommy Award will be tough to give out tonight.” Because we watched all the games, this was a part of the broadcast that intrigued us. Additionally, the Celtics' style of play makes it clear that effort is non-negotiable for their players. It makes us appreciate all players on all levels, whether they are scoring or drawing charges. We, as Celtics fans, appreciate the team for their meaningful contributions in any form.
\subsection{The Tommy Award}
The Tommy Award is named after 10x NBA champion and 39-year broadcaster for the Boston Celtics, Tommy Heinsohn. Over the course of Heinsohn’s career, he wanted to add something different to the broadcast. Instead of rewarding the best player every night for scoring the most points, he wanted to reward the players who did the “little things”. Getting steals, blocking shots, and drawing charges are examples of “little things” that can impact winning just as much as points. So NBC Sports Boston (the Celtics' home channel) decided to create an award in Heinsohn’s honor called the “Tommy Award”. The Tommy Award is given to one Celtics player after every game (win or loss) for demonstrating hustle and energy that helps their team win.

\subsection{Research Goal}

Although the Tommy Award is specific to the Celtics, the traits it rewards---hustle, defensive activity, rebounding, and low-usage impact---apply to players across the NBA. The award can tell us which players are impacting the game in ways that are not often discussed (deflections, charges drawn). As a result, these players are lesser known because their impact is overshadowed by their flashy teammates who score the most. So, by applying what the Tommy Award tells us to all NBA players, we can identify niche players whom NBA front offices (GMs, scouts, analysts) can add to their teams to ultimately help them win a championship.

The research paper will continue with the following sections. The Data Feature Construction section covers how we collected the data and what statistics were used in our model. The Methodology section will go in-depth on the models used and their performance metrics. Followed by the Results section, where the best-performing model was used to predict the Tommy Award for all other NBA teams. Finally, the Discussion of the results, followed by a Conclusion section which goes over what was learned, what can be improved, and next steps.

\section{Previous Work}

Chris Forsberg, whom I originally reached out to for the Tommy Award data, wrote an article documenting the history of the Tommy Award. In that article, there are tables showing the all-time leader in Tommy Award wins and the wins leader in each season dating back to the award's inception (2001-02). As far as showing who won each game and what their stats were, there was nothing, so that is where we decided to delve in.

As for previous work on measuring player performance in the NBA, ``Estimating an NBA Player's Impact on His Team's Chances of Winning'' by Deshpande and Jensen argues that player value should not be measured solely by traditional box score totals. Instead, the paper argues that context matters, since not every play has the same effect on the outcome of a game. A basket, rebound, or defensive play made in a high-leverage situation can matter much more than the same action in a less important moment. Because of this, the paper uses a win probability framework to evaluate how much a player's actions actually change his team's chances of winning.

Similarly, the Tommy Award aims to measure player value beyond the traditional box score, emphasizing hustle statistics. Timeliness is crucial for the Tommy Award because a player who is substituted in when they are not guaranteed playing time (such as subbing in for an injured player or getting more minutes than usual) has a hard time making an impact. If they do make a positive impact, they become a strong candidate for the award.

Additionally, ``Improved NBA Adjusted +/- Using Regularization and Out-of-Sample Testing'' by Joseph Sill focuses on a different but equally important issue in sports analytics: how to better estimate player impact while avoiding unstable or misleading results. The paper explains that adjusted plus-minus is useful because it attempts to isolate a player's contribution while accounting for teammates and opponents on the floor. However, the paper also shows that this kind of model can easily become noisy and overfit the training data. To address this, Sill applies regularization and evaluates the model using out-of-sample testing, meaning the model is judged by how well it performs on new data rather than just how well it fits old data. This is significant in a broader sports-prediction sense because it highlights that strong predictive models need to generalize beyond the training set and that balancing complexity with stability is a major part of effective sports analytics.

Although this project does not use adjusted plus-minus directly, Sill's work is relevant because it emphasizes the importance of testing models on unseen data. Since NBA player roles and team styles change over time, I used a chronological split rather than a random split. This allowed the model to train on older Celtics seasons and be evaluated on newer games, which better reflects the real-world task of predicting future Tommy Award winners.

\section{Data and Feature Construction}

\subsection{Tommy Award Data Collection}
The data for the Tommy Award winners was essential for this project. Originally, we tried manually collecting data from NBC Sports Boston's Instagram and website, where they posted the Tommy Award winner after each game. This only worked for the past two seasons (2024-25, 2025-26) and the 2016-17 season. Having three seasons' worth of data was simply not enough to do a thorough analysis; we would need as many seasons as possible. To expand the dataset, we contacted Chris Forsberg, a Celtics reporter for NBC Sports Boston, who directed me to Jason Levine, an executive producer at NBC Sports Boston. Levine provided Tommy Award data dating back to the award's inception. We asked for the 2016-17 season onward, as that's when the NBA started tracking hustle stats (deflections and charges drawn are the two notable hustle stats the NBA tracks, and we used them in this project). Again, hustle starts are an integral part of the Tommy Award, so the decision to start with when hustle stats were introduced in the NBA was an important one.

\subsection{NBA API Data Collection}
Once we had the Tommy Award data, retrieving per-game player stats from the NBA API was much easier. In total, we had 37 numerical features: 25 from the NBA API and the rest engineered by us. Of the 25 NBA API statistics, 15 of them are traditional box score stats. These include: minutes played, points, rebounds (offensive, defensive, total), assists, field goals made/attempted (for 3 pointers as well), free throws made, steals, blocks, and personal fouls. In addition to traditional box score stats, the NBA API provides analytics such as net rating, usage rate, deflections, charges drawn, and per-minute stats. Table~\ref{tab:nba-box} shows a full list of stats I used, including combinations of traditional stats like stocks and hustle proxy.

\begin{table}[htbp]
  \centering
  \caption{NBA-provided statistics used in the model.}
  \label{tab:nba-box}
  \begin{tabular}{@{}ll@{}}
    \toprule
    \textbf{Abbreviation} & \textbf{Statistic} \\
    \midrule
    MIN & Minutes Played \\
    PTS & Points \\
    AST & Assists \\
    OREB & Offensive Rebounds \\
    DREB & Defensive Rebounds \\
    REB & Total Rebounds \\
    STL & Steals \\
    BLK & Blocks \\
    TOV & Turnovers \\
    PF & Personal Fouls \\
    +/- & Plus/Minus \\
    FGM & Field Goals Made \\
    FGA & Field Goals Attempted \\
    FG3M & 3-Point Field Goals Made \\
    FG3A & 3-Point Field Goals Attempted \\
    FTM & Free Throws Made \\
    \bottomrule
  \end{tabular}
\end{table}

\subsection{Data Cleaning}
For each Tommy Award winner, I created player-level rows containing each player's stats for that game, using stats from the NBA API. This included the winner of the Tommy Award, as well as the other players in the game. Then, to use the Tommy Award data, I had to ensure that winners were marked as 1 in the CSV and losers as 0. I removed blank winners, winners that were not actual players (once a commentator won it, and multiple times, the ``stay ready'' group won it). Also, many of the winners were announced after midnight because the game was on the West Coast or had gone to overtime. We had to manually fix the dates for those games so the NBA API could correctly identify the games and add the desired stats to each player row. Fortunately, the Tommy Award data I received included the opponent for each game, so if the date was off, I could match the winner to the opponent with the official schedule to confirm the correct winner on the correct date.

\subsection{Engineered Features}
On top of the statistics the NBA provides, we wanted to create our own features to better value the team's unsung players. We did this by making features that prioritized high net rating and low usage. Net rating measures how much better or worse a team performs per 100 possessions while a player is on the floor. A positive net rating suggests that the team is outscoring opponents during that player's minutes. When predicting the Tommy Award, we want to reward players with consistently high net ratings. Especially if that player consistently takes fewer shots and has fewer touches (we call this usage), that is a sign they are contributing positively without using a large share of offensive possessions. Those players are doing the ``little things'' we mentioned earlier and that our data shows (deflections, charges drawn), which are helping their team win beyond just scoring points. By combining net rating and usage rate, I engineered a feature called impact efficiency. The equation for impact efficiency is
\[
  net\_rating / usage\_rate,
\]
and it rewards players who have a high net rating and a low usage rate. Remember, usage rate is the percentage of the team's plays that are ended by a particular player. A player with a low usage rate might sit in the corner or shoot 3s, or a big man who sets a lot of screens; either way, they are not their team's primary offensive option. Despite this, if their net rating is still high (their team goes up when they are on the floor), it means they are doing other things that aren't shooting/scoring to help their team win. This could be grabbing rebounds, getting steals, things that aren't scoring, but are still positively impacting their team. This player archetype is what we are seeking with the Tommy Award.

The second feature engineered statistic is role outperformance. Similar to impact efficiency, it involves both net rating and usage rate in the equation:
\[
  net\_rating * (1 - usage\_rate).
\]
Contrasting impact efficiency, this feature still rewards players with strong on-court impact but reduces the extreme effects that can occur when dividing by very small usage rates.

I also included rank stats, which sort teammates for the traditional statistical columns. So for a given game, Jayson Tatum may have 1 in the points column since he was first in points, and that same structure applies for the other statistics. The figure below shows the full list of my statistical columns.

\begin{table}[htbp]
  \centering
  \caption{Additional engineered features: impact efficiency, role outperformance, and rank statistics (column names as in \texttt{random\_forest\_classifier.ipynb}).}
  \label{tab:engineered}
  \small
  \begin{tabular}{@{}lp{8.8cm}@{}}
    \toprule
    \textbf{Feature} & \textbf{Role} \\
    \midrule
    \texttt{impact\_efficiency} & Net rating divided by usage rate (with numerical stabilization in code). \\
    \texttt{role\_outperformance} & Net rating $\times (1 - \texttt{usage\_rate})$. \\
    \texttt{stocks} & Steals + blocks. \\
    \texttt{hustle\_proxy} & (Offensive rebounds + steals + blocks) / minutes played \\
    Per-minute rates & \texttt{points\_per\_min}, \texttt{oreb\_per\_min}, \texttt{reb\_per\_min}, \texttt{ast\_per\_min}, \texttt{stocks\_per\_min}. \\
    Rank stats & Within-game teammate ranks such as \texttt{points\_rank}, \texttt{reboundsOffensive\_rank}, \texttt{reboundsTotal\_rank}, \texttt{assists\_rank}, \texttt{steals\_rank}, \texttt{blocks\_rank}, \texttt{plusMinusPoints\_rank}, \texttt{minutes\_decimal\_rank}, \texttt{stocks\_rank}, \texttt{hustle\_proxy\_rank}. \\
    \bottomrule
  \end{tabular}
\end{table}

\section{Methodology}

\subsection{Train/Validation/Test Split}
To create a successful model, we need a way to split our data into training and test sets, and chronological splitting is a way to do that. Instead of randomly dividing the dataset into folds, chronological splitting preserves the order of observations. Earlier seasons (2016-17 - 2022-23) are used as the training set, a later season (2023-24) as the validation set, and the most recent seasons (2024-25 - 2025-26) as the test set. Chronological splitting also helps prevent data leakage, since we do not use future games' outcomes to predict outcomes from earlier games. This approach is more appropriate than cross-validation for predicting the Tommy Award because player and league trends change over time, and chronological splitting reflects that.

\subsection{Models}
After assembling all features, we selected models for training. First, Ridge Logistic Regression served as a useful baseline because it allowed us to estimate the probability that each player would win while reducing the risk that noisy or correlated basketball statistics would dominate the model. Next, we implemented Lasso Logistic Regression, which can shrink coefficients to zero for effective feature selection. Finally, Elastic-net Logistic Regression combines elements of both Ridge and Lasso, maintaining Ridge's stability while selectively shrinking coefficients like Lasso. This approach effectively balances keeping relevant features with reducing the impact of less important ones.

Next, I explored tree-based models. First, I employed a Decision Tree, which splits data into branches based on feature values, illustrating how specific statistics can distinguish Tommy Award winners from non-winners. However, a single Decision Tree can be overly sensitive to the training data. To address this, I used a Random Forest, which creates multiple decision trees and combines their outputs for more robust and stable predictions. This approach reduces overfitting and captures more data patterns. Finally, I applied XGBoost, a sequential tree-building model in which each new tree corrects the previous ones' errors, making it effective at identifying complex patterns. Overall, these tree-based methods are crucial because they can model non-linear relationships and interactions between features that logistic regression might miss.

\subsection{Hyperparameter Tuning}
After creating our models, we want to tune their hyperparameters to achieve the best performance. One way to do this is, is to use Optuna. Optuna is a hyperparameter optimization tool that automatically tests different hyperparameter combinations and selects those that produce the best results. It helps identify the settings that yield the strongest model performance, making the tuning process more efficient than manually selecting values.

\subsection{Evaluation Metrics}
The two most crucial metrics for my project were top-3 accuracy and PR-AUC (Precision-Recall Area under the Curve). We chose top-3 because, in a game, even though there is one winner, two or three players typically perform well and are all likely to win. Thus, top-3 measures the percentage of times the actual winner is among the three highest-ranked predicted winners. Since my data is highly imbalanced (very few winners and many losers), PR-AUC indicates how effectively the model distinguishes actual Tommy Award winners from all non-winners. A higher PR-AUC score signifies better model identification of true winners.

I also included the Brier score and log loss as metrics. These are valuable because they evaluate the model's predicted probabilities, not just whether the predictions are correct. The Brier score measures how close the predicted probabilities are to the actual outcomes; thus, a lower score indicates more accurate probability estimates. Log loss is similar but penalizes the model more when it is confident and wrong, making it particularly useful for assessing not only the model's ability to identify Tommy Award winners but also its assignment of reasonable probabilities. For both metrics, lower values are better, indicating predicted probabilities that align more closely with the actual results.

I also used top-1 accuracy, which measures how often the model's highest-ranked player in a game was the actual Tommy Award winner. This is the strictest version of the prediction task. However, because the Tommy Award often has several reasonable candidates in a given game, top-3 accuracy is more informative. A model may still be useful if it consistently places the actual winner among the top few candidates, even when it does not rank that player first.

\section{Results}

\subsection{Model Comparison}
For my models, I used Lasso, Ridge, and Elastic Net Logistic Regression; Decision Tree, Random Forest, and XGBoost. I had five evaluation metrics: Top-1 accuracy, Top-3 accuracy, PR-AUC, Brier Score, and Log Loss. Below are the results

\begin{table}[htbp]
  \centering
  \caption{Model comparison on held-out seasons (values from project notebooks).}
  \label{tab:model-results}
  \begin{tabular}{@{}lrrrr@{}}
    \toprule
    \textbf{Model} & \textbf{Top-3 Acc.} & \textbf{PR-AUC} & \textbf{Brier Score} & \textbf{Log Loss} \\
    \midrule
    Random forest & 0.7853 & 0.3630 & 0.0709 & 0.2339 \\
    Elastic-net logistic & 0.7485 & 0.3143 & 0.1685 & 0.5024 \\
    Lasso logistic & 0.7423 & 0.3150 & 0.1695 & 0.5070 \\
    XGBoost & 0.7362 & 0.3363 & 0.1189 & 0.3514 \\
    Ridge logistic & 0.7362 & 0.3286 & 0.1774 & 0.5293 \\
    Decision tree & 0.6994 & 0.2383 & 0.0801 & 0.3945 \\
    \bottomrule
  \end{tabular}
\end{table}

Random Forest performed best across all five evaluation metrics, suggesting that it was better suited to the Tommy Award prediction task than the linear models, a single Decision Tree, or XGBoost. This result makes sense because the Tommy Award is not based on a single statistic. Instead, it depends on combinations of scoring, rebounding, defensive activity, usage, and team impact. Random Forest captured these interactions while avoiding the instability of a single Decision Tree. Although XGBoost is often a strong model, it may have been too sensitive to the noise in this dataset. In addition to the stat combinations, the Tommy Award is somewhat subjective because it can be awarded for a memorable hustle play, a late-game moment, or a broadcast narrative that does not fully appear in the box score. Since XGBoost builds trees sequentially to correct previous errors, it may have overfit to these smaller patterns. Random Forest was more stable because it built many independent trees and averaged their predictions, allowing it to capture broad feature interactions without overreacting to noisy individual examples.

\subsection{Best Model Performance}
After achieving 78\% top-3 and 42\% top-1 accuracy with Random Forest, I applied the same model to the remaining 29 NBA teams. I did this by taking the same Random Forest settings from optuna and, with the NBA API, gathered all 29 teams' 164 games across the 2024-25 and 2025-26 seasons. This process with Random Forest took approximately four hours. To compare players across the league, I calculated predicted Tommy Award wins per 60 minutes, requiring at least 300 minutes played. This prevented players with very small sample sizes from appearing artificially high and allowed lower-minute-role players to be compared more fairly with high-minute stars.

\subsection{NBA-Wide Tommy Award Predictions}

{\footnotesize
\setlength{\tabcolsep}{4pt}
\renewcommand{\arraystretch}{1.05}
\begin{longtable}{@{}r l l r r r@{}}
  \caption{Predicted Tommy Award wins per 60 minutes: top 100 player--team stints (at least 300 minutes played). Full ranking: \texttt{results\_per60\_table\_body.tex} (regenerate with \texttt{python3 scripts/build\_results\_per60\_table\_tex.py}).}\label{tab:per60}\\
  \toprule
  \textbf{\#} & \textbf{Player} & \textbf{Team} & \textbf{Wins} & \textbf{Min} & \textbf{Per-60} \\
  \endfirsthead
  \multicolumn{6}{c}{\textit{Table \ref{tab:per60} continued}}\\
  \toprule
  \textbf{\#} & \textbf{Player} & \textbf{Team} & \textbf{Wins} & \textbf{Min} & \textbf{Per-60} \\
  \endhead
  \midrule
  \multicolumn{6}{r}{\textit{Continued on next page}}\\
  \endfoot
  \bottomrule
  \endlastfoot
  % !TEX root = results_per60_table_standalone.tex
% This file is table *rows* only, not a full document. In TeXShop, Cmd-T builds the root above.
% On the command line: pdflatex results_per60_table_standalone.tex

        1 & Giannis Antetokounmpo & Milwaukee Bucks & 78 & 3327.9 & 1.4063 \\
        2 & Nikola Jokić & Denver Nuggets & 100 & 4835.7 & 1.2408 \\
        3 & Victor Wembanyama & San Antonio Spurs & 70 & 3392.9 & 1.2379 \\
        4 & Zion Williamson & New Orleans Pelicans & 47 & 2698.3 & 1.0451 \\
        5 & Mark Williams & Charlotte Hornets & 18 & 1171.1 & 0.9222 \\
        6 & Jalen Duren & Detroit Pistons & 60 & 4009.4 & 0.8979 \\
        7 & Ty Jerome & Memphis Grizzlies & 5 & 338.7 & 0.8858 \\
        8 & Luka Dončić & Los Angeles Lakers & 48 & 3272.7 & 0.8800 \\
        9 & Anthony Davis & Los Angeles Lakers & 20 & 1439.5 & 0.8336 \\
        10 & Shai Gilgeous-Alexander & Oklahoma City Thunder & 67 & 4856.6 & 0.8277 \\
        11 & Domantas Sabonis & Sacramento Kings & 40 & 2992.5 & 0.8020 \\
        12 & Jonas Valančiūnas & Washington Wizards & 13 & 982.4 & 0.7939 \\
        13 & Jimmy Butler III & Miami Heat & 10 & 765.8 & 0.7835 \\
        14 & Ivica Zubac & LA Clippers & 50 & 3952.7 & 0.7590 \\
        15 & Joel Embiid & Philadelphia 76ers & 22 & 1775.1 & 0.7436 \\
        16 & Anthony Davis & Dallas Mavericks & 11 & 891.6 & 0.7402 \\
        17 & Stephen Curry & Golden State Warriors & 43 & 3581.2 & 0.7204 \\
        18 & Karl-Anthony Towns & New York Knicks & 57 & 4838.5 & 0.7068 \\
        19 & Luka Dončić & Dallas Mavericks & 9 & 784.5 & 0.6883 \\
        20 & Cam Whitmore & Washington Wizards & 4 & 354.0 & 0.6779 \\
        21 & Alperen Sengun & Houston Rockets & 54 & 4792.8 & 0.6760 \\
        22 & Danté Exum & Dallas Mavericks & 4 & 372.3 & 0.6446 \\
        23 & Collin Sexton & Charlotte Hornets & 10 & 937.5 & 0.6400 \\
        24 & Marvin Bagley III & Dallas Mavericks & 5 & 470.2 & 0.6380 \\
        25 & Tyrese Maxey & Philadelphia 76ers & 49 & 4621.4 & 0.6362 \\
        26 & Mark Williams & Phoenix Suns & 15 & 1416.1 & 0.6356 \\
        27 & Chris Boucher & Toronto Raptors & 9 & 862.3 & 0.6262 \\
        28 & Nikola Vučević & Chicago Bulls & 39 & 3758.6 & 0.6226 \\
        29 & Devin Booker & Phoenix Suns & 51 & 4941.4 & 0.6193 \\
        30 & Jaren Jackson Jr. & Memphis Grizzlies & 37 & 3589.5 & 0.6185 \\
        31 & Jarrett Allen & Cleveland Cavaliers & 38 & 3815.3 & 0.5976 \\
        32 & Trae Young & Atlanta Hawks & 30 & 3019.2 & 0.5962 \\
        33 & Micah Potter & Indiana Pacers & 9 & 908.3 & 0.5945 \\
        34 & John Collins & Utah Jazz & 12 & 1220.5 & 0.5899 \\
        35 & Chet Holmgren & Oklahoma City Thunder & 28 & 2873.8 & 0.5846 \\
        36 & Josh Minott & Brooklyn Nets & 3 & 308.1 & 0.5843 \\
        37 & Kristaps Porziņģis & Atlanta Hawks & 4 & 413.1 & 0.5809 \\
        38 & Donovan Clingan & Portland Trail Blazers & 33 & 3417.4 & 0.5794 \\
        39 & Cam Thomas & Brooklyn Nets & 13 & 1363.1 & 0.5722 \\
        40 & Jakob Poeltl & Toronto Raptors & 27 & 2834.9 & 0.5715 \\
        41 & Cade Cunningham & Detroit Pistons & 44 & 4623.7 & 0.5710 \\
        42 & Walker Kessler & Utah Jazz & 18 & 1894.3 & 0.5701 \\
        43 & Daniel Gafford & Dallas Mavericks & 23 & 2420.6 & 0.5701 \\
        44 & Deandre Ayton & Portland Trail Blazers & 11 & 1205.6 & 0.5474 \\
        45 & Moritz Wagner & Orlando Magic & 9 & 990.8 & 0.5450 \\
        46 & Jordan Poole & Washington Wizards & 18 & 2000.9 & 0.5398 \\
        47 & Evan Mobley & Cleveland Cavaliers & 38 & 4240.9 & 0.5376 \\
        48 & Kel'el Ware & Miami Heat & 28 & 3126.3 & 0.5374 \\
        49 & Lauri Markkanen & Utah Jazz & 26 & 2918.5 & 0.5345 \\
        50 & Trendon Watford & Brooklyn Nets & 8 & 912.7 & 0.5259 \\
        51 & John Collins & LA Clippers & 16 & 1871.4 & 0.5130 \\
        52 & Jusuf Nurkić & Charlotte Hornets & 4 & 471.1 & 0.5095 \\
        53 & Tristan Vukcevic & Washington Wizards & 10 & 1183.5 & 0.5070 \\
        54 & Kristaps Porziņģis & Golden State Warriors & 3 & 355.8 & 0.5059 \\
        55 & Brandon Ingram & New Orleans Pelicans & 5 & 595.0 & 0.5042 \\
        56 & Rudy Gobert & Minnesota Timberwolves & 40 & 4768.6 & 0.5033 \\
        57 & D'Angelo Russell & Brooklyn Nets & 6 & 717.1 & 0.5020 \\
        58 & Jusuf Nurkić & Utah Jazz & 9 & 1083.0 & 0.4986 \\
        59 & Bennedict Mathurin & LA Clippers & 6 & 729.1 & 0.4938 \\
        60 & Zach Collins & Chicago Bulls & 6 & 735.8 & 0.4893 \\
        61 & Scottie Barnes & Toronto Raptors & 39 & 4814.9 & 0.4860 \\
        62 & Josh Giddey & Chicago Bulls & 31 & 3848.0 & 0.4834 \\
        63 & Jalen Brunson & New York Knicks & 39 & 4890.6 & 0.4785 \\
        64 & LeBron James & Los Angeles Lakers & 35 & 4433.5 & 0.4737 \\
        65 & Tyrese Haliburton & Indiana Pacers & 19 & 2450.4 & 0.4652 \\
        66 & Day'Ron Sharpe & Brooklyn Nets & 16 & 2067.4 & 0.4643 \\
        67 & Andrew Wiggins & Golden State Warriors & 10 & 1296.3 & 0.4629 \\
        68 & Alex Sarr & Washington Wizards & 24 & 3118.9 & 0.4617 \\
        69 & LaMelo Ball & Charlotte Hornets & 27 & 3522.0 & 0.4600 \\
        70 & Zach Edey & Memphis Grizzlies & 13 & 1699.6 & 0.4589 \\
        71 & Pascal Siakam & Indiana Pacers & 35 & 4604.5 & 0.4561 \\
        72 & Julian Reese & Washington Wizards & 3 & 401.5 & 0.4483 \\
        73 & Anthony Edwards & Minnesota Timberwolves & 37 & 5007.6 & 0.4433 \\
        74 & Coby White & Charlotte Hornets & 3 & 406.4 & 0.4429 \\
        75 & Paolo Banchero & Orlando Magic & 30 & 4083.2 & 0.4408 \\
        76 & Cole Anthony & Orlando Magic & 9 & 1234.1 & 0.4375 \\
        77 & Jalen Green & Phoenix Suns & 6 & 828.3 & 0.4346 \\
        78 & Trey Murphy III & New Orleans Pelicans & 30 & 4194.9 & 0.4291 \\
        79 & Kevin Durant & Phoenix Suns & 16 & 2264.7 & 0.4239 \\
        80 & Orlando Robinson & Toronto Raptors & 5 & 712.7 & 0.4209 \\
        81 & Jimmy Butler III & Golden State Warriors & 15 & 2162.0 & 0.4163 \\
        82 & Desmond Bane & Orlando Magic & 19 & 2756.1 & 0.4136 \\
        83 & Precious Achiuwa & Sacramento Kings & 12 & 1745.2 & 0.4126 \\
        84 & Collin Sexton & Utah Jazz & 12 & 1757.8 & 0.4096 \\
        85 & Nick Richards & Charlotte Hornets & 3 & 440.3 & 0.4088 \\
        86 & Jock Landale & Atlanta Hawks & 3 & 446.0 & 0.4036 \\
        87 & Ja Morant & Memphis Grizzlies & 14 & 2087.6 & 0.4024 \\
        88 & Marvin Bagley III & Washington Wizards & 6 & 897.2 & 0.4012 \\
        89 & James Harden & LA Clippers & 29 & 4348.3 & 0.4002 \\
        90 & Jalen Hood-Schifino & Philadelphia 76ers & 2 & 300.6 & 0.3993 \\
        91 & Tyler Herro & Miami Heat & 25 & 3757.8 & 0.3992 \\
        92 & Robert Williams III & Portland Trail Blazers & 9 & 1359.6 & 0.3972 \\
        93 & Michael Porter Jr. & Brooklyn Nets & 11 & 1688.6 & 0.3908 \\
        94 & Dillon Brooks & Phoenix Suns & 11 & 1701.6 & 0.3879 \\
        95 & Taylor Hendricks & Memphis Grizzlies & 4 & 626.3 & 0.3832 \\
        96 & Joan Beringer & Minnesota Timberwolves & 2 & 314.1 & 0.3820 \\
        97 & Kevin Durant & Houston Rockets & 18 & 2839.8 & 0.3803 \\
        98 & Jalen Johnson & Atlanta Hawks & 24 & 3815.8 & 0.3774 \\
        99 & Cedric Coward & Memphis Grizzlies & 10 & 1597.9 & 0.3755 \\
        100 & Bam Adebayo & Miami Heat & 31 & 5038.4 & 0.3692 \\
\end{longtable}
}

\subsection{Feature Importance}
Table~\ref{tab:fi} lists the top fifteen random forest feature importances; values match the bar chart in the analysis notebook (three decimal places as displayed there). The full ranking of all thirty-nine numeric features is in \texttt{results\_rf\_feature\_importance\_body.tex}, regenerated by \texttt{python3 scripts/export\_rf\_feature\_importances\_tex.py}.

{\footnotesize
\setlength{\tabcolsep}{5pt}
\renewcommand{\arraystretch}{1.1}
\begin{table}[htbp]
  \centering
  \caption{Random Forest: top fifteen feature importances.}\label{tab:fi}
  \begin{tabular}{@{}r l r@{}}
    \toprule
    \textbf{Rank} & \textbf{Feature} & \textbf{Importance} \\
    \midrule
    1 & \texttt{points\_per\_min} & 0.088 \\
    2 & \texttt{points} & 0.075 \\
    3 & \texttt{fieldGoalsMade} & 0.069 \\
    4 & \texttt{points\_rank} & 0.060 \\
    5 & \texttt{usage\_rate} & 0.045 \\
    6 & \texttt{fieldGoalsAttempted} & 0.043 \\
    7 & \texttt{minutes\_decimal} & 0.039 \\
    8 & \texttt{reb\_per\_min} & 0.034 \\
    9 & \texttt{threePointersMade} & 0.033 \\
    10 & \texttt{impact\_efficiency} & 0.032 \\
    11 & \texttt{ast\_per\_min} & 0.031 \\
    12 & \texttt{role\_outperformance} & 0.029 \\
    13 & \texttt{plusMinusPoints\_rank} & 0.027 \\
    14 & \texttt{hustle\_proxy} & 0.027 \\
    15 & \texttt{stocks\_per\_min} & 0.026 \\
    \bottomrule
  \end{tabular}
\end{table}
}

\section{Discussion}

Going into modeling, we selected hustle proxy, charges drawn, deflections, and total rebounds because, this season, watching Celtics games, those were the things that stood out as really important for role players to do to win the Tommy Award. These are also generally important things to help your team win, so doing these at a high level helps players win more Tommy Awards.

Ultimately, points and field goals made were among the top 3 most important statistics for predicting the Tommy Award. This result complicates the original assumption behind the Tommy Award. Although the award is meant to recognize hustle and energy, the model suggests that traditional scoring statistics still play a major role. This may be because high-scoring players are more visible, but it may also be because the best players often contribute in many areas at once. In other words, the Tommy Award does not ignore scoring as much as expected; instead, it appears to reward players who combine scoring with other forms of impact. However, MVPs like Nikola Joki\'c and Giannis Antetokounmpo were among the league leaders, while lesser-known players like Ty Jerome and Mark Williams were in the top 10. Additionally, many centers/big men were in the top 50. This is because we put a lot of value on offensive rebounding. That is because this season the Celtics made it an emphasis for all players. However, for all teams, it will obviously be easier for bigs to get offensive rebounds and rebounds in general.

\section{Conclusion and Future Work}

We found that points, field goals made, and field goals attempted were among the most influential statistics for the Tommy Award. In addition, Random Forest achieved 78\% top-3 accuracy and 42\% top-1 accuracy. We were able to apply this to the entire NBA. Players such as Ty Jerome, Mark Williams, and Cam Whitmore were among the most important takeaways from the project because they represent the type of overlooked contributors the Tommy Award is designed to identify. These are players who may not always receive national attention, but the model suggests that they contribute in ways that align with the Tommy Award profile. These are the players that NBA front offices need to look out for as potential pieces for a championship-contending team.

The Tommy Award is useful because it captures a broader type of player value than traditional scoring-focused awards. By training a model on Celtics Tommy Award winners, we can identify NBA players who fit that same profile across the league. However, the model also shows that traditional scoring still matters more than expected, meaning the Tommy Award is not purely a hustle award but a combination of visibility, production, and role-based impact.

Future work should address the class imbalance in the dataset. Since each game has only one winner and many non-winners, the model is trained on far more negative examples than positive ones. One possible solution is to use resampling methods such as oversampling winners, undersampling non-winners, or synthetic sampling techniques. However, any synthetic data would need to be carefully created to avoid producing unrealistic basketball statistics.

Furthermore, I want to improve our feature engineering and explore how to compute statistics for different types of players. Because we valued offensive rebounds in the project, we could make offensive rebounds by guards and wings more valuable than those by centers. Additionally, we could add more features from the NBA stats dashboard to our list, such as loose balls recovered, contested shots, and box-outs.

Finally, we want to try to work with play-by-play data. We want to see which sequence in the game made a certain player very likely to win the Tommy Award. A series of key defense stops, big threes made, and diving on the floor for a loose ball. What play contributed to that player winning the award and or to their team winning the game?

\section{References}
\begin{enumerate}
  \item Deshpande, S. K., \& Jensen, S. T. (2016). Estimating an NBA player's impact on his team's chances of winning. \textit{Journal of Quantitative Analysis in Sports, 12}(2), 51--72. \url{https://doi.org/10.1515/jqas-2015-0027}

  \item Forsberg, C. (2020, May 6). \textit{How the Tommy Award became a measuring stick for Celtics players}. NBC Sports Boston. \url{https://www.nbcsportsboston.com/nba/boston-celtics/how-the-tommy-award-became-a-measuring-stick-for-celtics-players/336968/}

  \item Houde, M. (2021). \textit{Predicting the outcome of NBA games} [Honors thesis, Bryant University]. Bryant Digital Repository. \url{https://digitalcommons.bryant.edu/honors_data_science/1/}

  \item James, G., Witten, D., Hastie, T., Tibshirani, R., \& Taylor, J. E. (2023). \textit{An introduction to statistical learning: With applications in Python}. Springer. \url{https://doi.org/10.1007/978-3-031-38747-0}

  \item Kuehn, J. (2016). \textit{Accounting for complementary skill sets when evaluating NBA players' values to a specific team}. MIT Sloan Sports Analytics Conference. \url{https://cdn.prod.website-files.com/68d6be744d7efccc2207f571/68d6be744d7efccc22080648_Accounting%20for%20Complementary%20Skill%20Sets%20when%20Evaluating%20NBA%20Players%27%20Values%20to%20a%20Specific%20Team.pdf}

  \item NBC Sports Boston. (n.d.). \textit{The Tommy Award}. \url{https://www.nbcsportsboston.com/the-tommy-award/}

  \item Sarlis, V., Gerakas, D., \& Tjortjis, C. (2024). A data science and sports analytics approach to decode clutch dynamics in the last minutes of NBA games. \textit{Machine Learning and Knowledge Extraction, 6}(3), 2074--2095. \url{https://doi.org/10.3390/make6030102}

  \item Sill, J. (2010). \textit{Improved NBA adjusted +/- using regularization and out-of-sample testing}. MIT Sloan Sports Analytics Conference. \url{https://www.sloansportsconference.com/research-papers/improved-nba-adjusted-using-regularization-and-out-of-sample-testing}

  \item Williams, B., Schliep, E. M., Fosdick, B., \& Elmore, R. (2024). \textit{Expected points above average: A novel NBA player metric based on Bayesian hierarchical modeling}. arXiv. \url{https://doi.org/10.48550/arXiv.2405.10453}
\end{enumerate}

\end{document}
